\section{Naive Set theory}

\exercise Russell's Paradox shows that na\"{\i}ve set theory's assumption of the unrestricted comprehension principle -- for any property $P$, the set $\{x: P(x)\}$ exists -- leads to a contradiction. Consider the property $P(x): x \in x$. The unrestricted comprehension prinicple tells us that the set $
S = \{x : \sim P(x)\}$ exists. This is simply not possible, since we get $S \in S \iff S\not\in S$.\newline
To see this, note that if $S$ contains itself, then $S$ does not satisfy $\sim P(S)$, so $S\not\in S$, but now $S$ does satisfy $\sim P(S)$, so $S\in S$. i.e, $S\in S\iff S\not\in S$, and we know that only one of those alternatives is possible. A more detailed discussion can be found at \url{https://plato.stanford.edu/archives/fall2025/entries/russell-paradox/}
\exercise First, we check that the elements of $\mathscr{P}_{\sim}$ are nonempty. Indeed, since $\sim$ is an equivalence relation, we have that for all $a\in S: a\sim a$, so that $a\in [a]$ i.e $[a]$ is nonempty. Now we check that distinct elements of $\mathscr{P}_{\sim}$ are disjoint. Indeed, let $[a]\neq [b]$ be two elements of $\mathscr{P}_{\sim}$. If they did intersect at some $c$, we would have $a\sim c$ and $b\sim c$ which by transitivity would give $a\sim b$. This is simply not possible, because this implies by transitivity that for every $x\in S$ such that  $x\sim b$, we have $x\sim a$. i.e $[b]\subseteq [a]$. Symmetry gives $[a]\subseteq [b]$ i.e $[a]=[b]$, which contradicts our initial assumption of $[a]\neq [b]$. Ofcourse the union of the elements of $\mathscr{P}_{\sim}$ is $S$ because the elements of $\mathscr{P}_{\sim}$ are $[a]\supseteq \{a\}$ for each $a\in S$. This gives us \[\bigcup_{a\in S}[a] \supseteq S.\] The other inclusion is clear, since each $[a]$ is a subset of $S$. Thus we have $\bigcup_{x\in \mathscr{P}_{\sim}}x = \bigcup_{a\in S}[a] = S$ as desired. i.e, $\mathscr{P}_{\sim}$ is a partition of $S$.

\exercise let $\mathscr P = \{S_{\alpha} : \alpha\in I\}$ be a partition of $S$. We claim that the equivalence relation $\sim$ given by $a\sim b \iff \left(\exists \alpha\in I : a\in S_\alpha \text{ and } b\in S_\alpha\right)$. I.e two elements are related iff they lie in the same part of the partition. Indeed, this is an equivalence relation. For any $a\in S$, we have $a\sim a$ because $a\in S_\alpha$ for some $\alpha$ so $\sim$ is reflexive. Now suppose that $a\sim b$ and $b\sim c$. Since $b$ lies in exactly one $S_\alpha$, we have that $a$ and $c$ lie in the same $S_\alpha$, i.e $a\sim c$ so $\sim$ is transitive. Symmetry is clear, if $a$ lies in the same part as $b$, $b$ lies in the same part as $a$.

\exercise From what we have seen, it suffices to count the number of partitions of $\{1,2,3\}$, which are $\{\{1\},\{2\},\{3\}\}, \{\{1,2\},\{3\}\}, \{\{1,3\},\{2\}\}, \{\{2,3\},\{1\}\}, \{\{1,2,3\}\}$. I.e, there are $5$ equivalence relations on the set $\{1,2,3\}$

\exercise On the set $\{1,2,3\}$ define the following equivalence relation: \[1\sim 1, 2\sim 2, 3\sim 3, 3\sim 2, 2\sim 1, 2\sim 3, 1\sim 2.\] This is reflexive and symmetric but not transitive. Indeed, $1\sim 2$ and $2\sim 3$ but $1\not\sim 3$. If we use a relation that is only reflexive and symmetric but not transitive, then we will not get a partition. Indeed, the parts may intersect. For example, in the above relation we defined, if we let $[a] := \{b: a\sim b\}$, then $[1]\cap [2] = \{2\}$ even though $[1]\neq [2]$.

\exercise The relation $\sim$ is reflexive since indeed, $a - a = 0\in \mathbb{Z}$ for all $a\in \mathbb{R}$. It is symmetric, since the negative of an integer is an integer. Indeed, $a\sim b \implies a-b = k_1 \implies b-a = -k_1$. It is transitive, since the sum of two integers is an integer. Indeed, if $a\sim b$, $b\sim c$, we have $a-b = k_1, b-c = k_2$ for some $k_1,k_2\in \mathbb{Z}$ which gives $a - c = k_1 + k_2\in\mathbb{Z}$. Hence, it is an equivalence relation. In exactly the same way, we get that the relation $(a_1,a_2) \approx (b_1,b_2)\iff a_1 - b_1\in\mathbb{Z}\text{ and }a_2 - b_2\in \mathbb{Z}$ is an equivalence relation.\newline
As for finding ``compelling" descriptions, first relation can be seen geometrically as a `circle' in the sense that the set $\mathbb{R}/\sim$ can be identitfied with the set $\{e^{2\pi i x}\}$ via $[x]\mapsto e^{2\pi i x}$. It is important to note that since we are using the representative of the classes in defining our correspondence, that it is well defined. Similarly, the set $\mathbb{R}\times\mathbb{R}/\approx$ can be identitifed with a torus, via the same exponential map: $[(x,y)] \mapsto (e^{2\pi i x}, e^{2\pi i y})$ which one can again check is well defined.

